% DRAFT 0, 2026-11-11. Not compiled: no TeX installation was available when
% this was written. Not yet in the JAIR style: replace the preamble with the
% JAIR template and append its reproducibility checklist before submission.
%
% Every table in tables/ and every data file in data/ is generated by
% ../make_paper_assets.py from ../results/*.csv. Do not edit them by hand.
%
% Status of the text: first draft produced with AI assistance from the
% project's design notes. Under JAIR policy the author is responsible for
% every claim, proof and citation. Proofs have NOT been checked by a human.
% Entries in refs.bib marked "verified = {no}" have not been checked against
% the source.
\documentclass[11pt]{article}
\usepackage[margin=1in]{geometry}
\usepackage{amsmath,amssymb,amsthm}
\usepackage{booktabs}
\usepackage{natbib}
\usepackage{pgfplots}
\pgfplotsset{compat=1.17}
\usepackage{hyperref}

\newtheorem{proposition}{Proposition}
\newtheorem{lemma}{Lemma}
\newtheorem{corollary}{Corollary}
\theoremstyle{remark}
\newtheorem{remark}{Remark}

\newcommand{\E}{\mathbb{E}}
\newcommand{\Prob}{\mathbb{P}}
\newcommand{\ind}{\mathbf{1}}
\newcommand{\dlo}{d_{\mathrm{lo}}}
\newcommand{\dhi}{d_{\mathrm{hi}}}
\newcommand{\todo}[1]{\textbf{[TODO: #1]}}

\title{When Does Causal Correction Help Discount Recommendation?\\
Confounded Discounts in Negotiated Pricing}
\author{[Author name]\\ {[Affiliation]}}
\date{Draft 0, November 2026}

\begin{document}
\maketitle

\begin{abstract}
\textbf{Background.} Sellers of negotiated products increasingly learn a
win-probability model from past deals and use it to recommend discounts. The
salespeople who set those past discounts knew things about each deal that the
model does not, so discount is confounded with outcome.
\textbf{Objectives.} We ask when this confounding changes the recommended
discount, whether the common remedy of constraining the model to be monotone
in discount repairs the decision, and when an instrument built from
salesperson leniency helps or harms.
\textbf{Methods.} We analyse a single-index model of rep-driven confounding
and build a semi-synthetic benchmark with a known causal response, scored by
revenue regret, with controls for the strength and direction of confounding,
sample size, model capacity, and the validity of the instrument.
\textbf{Results.} Mild confounding barely moves the recommendation; strong
confounding collapses it, and which direction is harmful depends on the win
model. A monotone constraint repairs the fitted curve but not the decision,
and under a saturating response it lowers the recommendation and raises
regret. A leniency control function is close to an oracle on large books when
salespeople are assigned as if at random, and is worse than no correction when
lenient salespeople differ in skill or receive systematically different deals.
We prove that the first of these failures is undetectable from the deal data
and show that standard balance and first-stage checks pass under the second.
\textbf{Conclusions.} Causal correction for discount guidance is worth its
variance only when assignment is rule-based and the book is large. A
monotonicity belief is harmful as a constraint on an uncorrected model and
useful as a falsification check on a corrected one.
\end{abstract}

\section{Introduction}

A growing class of decision-support systems is trained on the logged decisions
of the people it is meant to advise. A model observes what experts did and
what happened, and then recommends what to do next. When the experts acted on
information the model cannot see, a model that predicts outcomes well can
still recommend badly, because the action and the outcome share a hidden
cause.

Negotiated pricing is a clean instance. In business-to-business sales a
representative (rep) chooses a discount for each deal. Vendors of pricing
software fit a model of the probability of winning as a function of deal
attributes and the discount, and recommend the discount that maximizes
expected revenue. Reps discount with the deal in front of them in mind: a
struggling deal may get a deeper discount to rescue it, or a strategically
important one may get a deeper discount to make sure of it. Either way the
discount carries information about the deal's prospects that is not recorded.

Practitioners know this. A common response is to constrain the win model to
be non-decreasing in discount, on the argument that a larger discount can
never hurt the chance of winning, so a fitted curve that slopes down must be
an artefact of confounding. The constraint is cheap, it is supported by
standard gradient-boosting libraries, and it produces curves that are easy to
defend to a customer. This paper started from a simple question about that
practice: does forcing the curve to have the right sign fix the
recommendation?

It does not. We show three things.

\begin{enumerate}
\item \textbf{The constraint repairs the picture, not the decision.} In a
linear-index model the monotone-constrained and unconstrained models recommend
the same discount for every strength and direction of confounding
(Corollary~\ref{cor:neutral}). When the true response saturates, the
constrained model recommends weakly less discount than the unconstrained one
and has weakly higher regret whenever the unconstrained model is already
under-discounting (Proposition~\ref{prop:saturating}). In our benchmark this
costs 4 to 11 points of revenue under moderate to strong confounding at the
capacity of a production model.
\item \textbf{Whether confounding matters depends on the model as much as on
the data.} Regret is flat in the bias for small biases and then rises sharply
(Proposition~\ref{prop:regret}). The direction of confounding that is harmful
flips between a model that under-fits the response and one that fits it well.
At realistic sample sizes, the capacity of the win model moves regret by more
than confounding does.
\item \textbf{An instrument from rep leniency works under conditions that
sales organizations rarely guarantee, and its failures are silent.} A control
function built on leave-one-out rep leniency is within a point of an oracle
on books of 20{,}000 deals with randomly assigned reps. If lenient reps are
also weaker or stronger sellers, the corrected model is worse than the
uncorrected one, and we prove that no test on the deal data can detect it
(Proposition~\ref{prop:undetectable}). If managers route deals to reps on
unrecorded deal quality, the same happens, and the balance test and
first-stage diagnostics that the examiner-design literature recommends pass.
\end{enumerate}

The contribution is an analysis and a benchmark, not a new estimator. The
analysis says when correction changes the decision. The benchmark, which we
release, lets methods for decision-focused learning and causal policy learning
be tested under controlled violations of their assumptions, with regret
measured against a known truth. Pricing is the worked setting; the questions
arise wherever a shape constraint is used to make a model fitted to logged
human decisions look causal.

\section{Setting and Related Work}
\label{sec:related}

\paragraph{Customized pricing.} In customized or negotiated pricing each
customer receives an individual quote and accepts or declines
\citep{phillips2015field}. The standard tool is a bid-response model, the
probability of acceptance as a function of price. \citet{phillips2015field}
study an indirect auto lender whose dealers negotiate rates, correct for price
endogeneity with a control function \citep{petrin2010control} using rates
offered for similar loans in other regions as the instrument, and report that
the correction raised their estimated price elasticities by 126\% to 285\%. \citet{karlinsky2024automating} model the pricing of individual
salespeople at a business-to-business retailer, instrument price with product
cost, and route quotes between model and salesperson. Neither uses the
salesperson as the instrument, and neither studies shape constraints.
\citet{bockelrickermann2023causal} benchmark estimators of bid-response
functions on semi-synthetic mortgage data under selection on observed
covariates; our benchmark differs in having unobserved confounding, a rep
structure, and regret in revenue terms. \todo{Read all three in full and
check each characterization; add Miao et al.\ and Gupta and Kamble after
reading.}

\paragraph{Examiner designs.} Using the leniency of a quasi-randomly assigned
decision-maker as an instrument is standard in economics
\citep{kling2006incarceration,dobbie2018effects}. \citet{frandsen2023judging}
give a joint test of the exclusion and monotonicity assumptions, and
\citet{chan2022selection} show that such designs fail when decision-makers
differ in skill and not only in threshold. Our
Proposition~\ref{prop:undetectable} is the form that failure takes in pricing.
\todo{Read Chyn, Frandsen and Leslie (2025) and cite.}

\paragraph{Decisions versus predictions.} Decision-focused learning trains
models for the quality of the downstream decision \citep{elmachtoub2022smart,
mandi2024decision}, assuming the training data sample the decision environment
faithfully. \citet{fernandezloria2022causal} argue that good decisions do not
require accurate effect estimates and that confounded data can still support
them. Our results agree for mild confounding and locate where the argument
stops holding.

\paragraph{Robustness and deferral under hidden confounding.}
\citet{kallus2021minimax} learn policies that minimize worst-case regret over
a sensitivity model, and \citet{conley2012plausibly} relax the exclusion
restriction to a bounded direct effect. \citet{ghoummaid2024act} learn when to
defer to a human expert under bounded hidden confounding. We use a
sensitivity range in the same spirit and find that choosing robustly dominates
deferring to a default in our setting.

\paragraph{Shape constraints.} Monotonicity is known to stabilize
nonparametric instrumental-variable estimation \citep{chetverikov2017nonparametric}.
We are not aware of an analysis of what a monotone constraint does to the
decision when no instrument is used. \todo{Search again before submission.}

\section{A Model of Rep-Driven Confounding}
\label{sec:theory}

We consider one cohort of deals; everything below holds conditional on
observed covariates. A deal has latent quality $q\sim N(0,1)$, seen by the rep
and not by the model. The rep sets the discount
\begin{equation}
d = m + \gamma q + \eta, \qquad \eta\sim N(0,\sigma^2)\ \text{independent of } q,
\label{eq:first}
\end{equation}
and the deal is won if
\begin{equation}
Y = \ind\{\alpha + g(d) + \beta q + \varepsilon > 0\}, \qquad
\varepsilon\sim N(0,1),
\label{eq:outcome}
\end{equation}
with $g$ increasing and $\beta>0$. Write $\sigma_d^2=\gamma^2+\sigma^2$ and
$\rho=\gamma/\sigma_d$ for the correlation between discount and quality. We
call $\rho<0$ \emph{rescue} confounding (deeper discounts on weaker deals) and
$\rho>0$ \emph{whale} confounding (deeper discounts on stronger deals). The
causal win curve is
\begin{equation}
p^*(d) = \Prob(Y=1\mid do(d)) = \Phi\!\left(\frac{\alpha+g(d)}{\sqrt{1+\beta^2}}\right).
\end{equation}
The seller chooses $d\in[\dlo,\dhi]\subset[0,1)$ to maximize expected revenue
per unit of list price, $R(d)=(1-d)\,p^*(d)$, with maximizer $d^*$. A model
that believes a curve $\hat p$ recommends
$\hat d=\arg\max_d (1-d)\hat p(d)$, taking the smallest maximizer, and its
regret is $R(d^*)-R(\hat d)$.

\subsection{What a predictive model learns}

\begin{proposition}[Observational curve]
\label{prop:obs}
Under \eqref{eq:first} and \eqref{eq:outcome},
\begin{equation}
\pi(d) := \Prob(Y=1\mid d) =
\Phi\!\left(\frac{\alpha + g(d) + \beta\rho\,(d-m)/\sigma_d}
{\sqrt{1+\beta^2(1-\rho^2)}}\right).
\label{eq:obs}
\end{equation}
If $g(d)=\theta d$ with $\theta>0$, then $\pi$ is a probit in $d$ with slope
$b_{\mathrm{obs}}=(\theta+\beta\rho/\sigma_d)/\sqrt{1+\beta^2(1-\rho^2)}$,
and: (i) $b_{\mathrm{obs}}>b^*=\theta/\sqrt{1+\beta^2}$ whenever $\rho>0$;
(ii) $b_{\mathrm{obs}}<0$ if and only if $\rho<-\theta\sigma_d/\beta$;
(iii) the confounding term vanishes at $d=m$, so the observational curve is
the causal curve rotated about the mean historical discount, up to scale.
\end{proposition}
\begin{proof}
$(q,d)$ is bivariate normal, so $q\mid d\sim N(\rho(d-m)/\sigma_d,\,1-\rho^2)$.
Write $q=\rho(d-m)/\sigma_d+\sqrt{1-\rho^2}\,\xi$ with $\xi\sim N(0,1)$
independent of $d$. Then $Y=1$ if and only if
$\alpha+g(d)+\beta\rho(d-m)/\sigma_d+\beta\sqrt{1-\rho^2}\,\xi+\varepsilon>0$,
and $\beta\sqrt{1-\rho^2}\,\xi+\varepsilon\sim N(0,1+\beta^2(1-\rho^2))$.
Parts (i) to (iii) follow by inspection; for (i) both the numerator and the
reciprocal of the denominator exceed their values at $\rho=0$.
\end{proof}

\begin{remark}
Rescue confounding does not always flatten the curve on the index scale: the
denominator in \eqref{eq:obs} shrinks as $|\rho|$ grows, and when
$\theta\sigma_d\gg\beta$ this can outweigh the numerator. To first order in
$\rho$ the bias has the sign of $\rho$.
\end{remark}

\subsection{What a monotone constraint does}

Let $\tilde\pi$ be the population log-loss minimizer over functions of $d$
that are non-decreasing, that is, the isotonic regression of $\pi$ with
respect to the distribution of $d$ \citep{robertson1988order}.

\begin{proposition}[Monotone fit, linear index]
\label{prop:mono}
Let $g(d)=\theta d$. If $b_{\mathrm{obs}}\ge 0$ then $\tilde\pi=\pi$. If
$b_{\mathrm{obs}}<0$ then $\tilde\pi$ is the constant $\Prob(Y=1)$.
\end{proposition}
\begin{proof}
A non-decreasing function is its own isotonic regression. The isotonic
regression of a strictly decreasing function is its mean with respect to the
weighting measure.
\end{proof}

\begin{corollary}[Decision neutrality]
\label{cor:neutral}
In the linear-index model the constrained and unconstrained models recommend
the same discount for every $\rho$.
\end{corollary}
\begin{proof}
If $b_{\mathrm{obs}}\ge0$ the two models coincide. If $b_{\mathrm{obs}}<0$,
then under the unconstrained belief $\pi$ is decreasing and under the
constrained belief $\tilde\pi$ is constant; in both cases $(1-d)\hat p(d)$ is
strictly decreasing and both recommend $\dlo$.
\end{proof}

So in the linear model the constraint changes what the user is shown, from a
curve on which discounts hurt to one on which they do nothing, and leaves the
recommendation and the regret unchanged. With a saturating response it does
worse.

\begin{proposition}[Saturating response]
\label{prop:saturating}
Let $g$ be increasing and concave. Then:
\begin{enumerate}
\item $\pi$ is increasing if $\rho\ge0$. If $\rho<0$ it is single-peaked: it
increases while $g'(d)>\beta|\rho|/\sigma_d$ and decreases afterwards.
\item If $\pi$ is single-peaked with peak at $d_p$, there is $t\le d_p$ such
that $\tilde\pi=\pi$ on $[\dlo,t]$ and $\tilde\pi$ is constant on $(t,\dhi]$.
\item Writing $d_u$ and $d_c$ for the unconstrained and constrained
recommendations, $d_c\le d_u$, and $d_c\le t$ in the single-peaked case. If
$\rho\ge0$ then $d_c=d_u$.
\item If $d_u\le d^*$ then $R(d^*)-R(d_c)\ge R(d^*)-R(d_u)$, strictly if
$d_c<d_u$.
\end{enumerate}
\end{proposition}
\begin{proof}[Proof sketch]
(1) The argument of $\Phi$ in \eqref{eq:obs} has derivative proportional to
$g'(d)+\beta\rho/\sigma_d$, and $g'$ is non-increasing. (2) This is the
pool-adjacent-violators solution for a function that rises and then falls: the
falling side is pooled with as much of the rising side as is needed for the
pooled mean to equal the value of the curve at the start of the pool
\citep{robertson1988order}. \todo{Give the argument in full, including the
boundary case $t=\dlo$.} (3) On $(t,\dhi]$, $(1-d)\tilde\pi(d)$ is strictly
decreasing, so $d_c\le t$. On $[\dlo,t]$ the two objectives coincide. If
$d_u\le t$ then $d_c=d_u$; otherwise $d_c\le t<d_u$. (4) $\Phi$ is
log-concave and $g$ is concave, so $\log R$ is concave on $[0,1)$ and $R$ is
increasing on $[\dlo,d^*]$.
\end{proof}

Part 4 has a condition. Lemma~\ref{lem:direction} below shows that mild rescue
confounding can push $d_u$ above $d^*$, and then a lower recommendation may
help. Part 3 is specific to a single-peaked curve; with a dip between two high
regions the constrained recommendation can be the higher one.
Figure~\ref{fig:curves} illustrates the mechanism: the constrained model
recommends the point at which its curve goes flat. The constraint converts
``the data say discount stops helping here'' into a hard stop, and under
rescue confounding the data say so too early.

\begin{figure}[t]
\centering
\begin{tikzpicture}
\begin{axis}[width=0.72\textwidth, height=0.42\textwidth,
  xlabel={Discount (\%)}, ylabel={Win probability},
  xmin=0, xmax=25, ymin=0.25, ymax=0.75,
  legend pos=south east, legend cell align=left]
\addplot[black, thick] table[x=d, y=causal] {data/curves.dat};
\addlegendentry{Causal}
\addplot[black, thick, dashed] table[x=d, y=obs] {data/curves.dat};
\addlegendentry{Observational}
\addplot[black, thick, dotted] table[x=d, y=mono] {data/curves.dat};
\addlegendentry{Monotone fit}
\end{axis}
\end{tikzpicture}
\caption{Win curves under rescue confounding ($\rho=-0.5$, $\alpha=-0.5$,
$\beta=0.55$, mean historical discount 12\%) with the saturating response used
in the benchmark. The revenue-maximizing discount is 8.9\% on the causal
curve, 5.7\% on the observational curve, and 3.1\% on the monotone fit, which
is where the fit goes flat.}
\label{fig:curves}
\end{figure}

\subsection{When the bias changes the decision}

\begin{proposition}[Regret]
\label{prop:regret}
Let $g$ be increasing and concave and let a model believe
$\hat p(d)=\Phi(\hat a+\hat b d)$.
\begin{enumerate}
\item $R$ is single-peaked, so regret is at most
$\max\{R(d^*)-R(\dlo),\,R(d^*)-R(\dhi)\}$.
\item If $d^*$ is interior, regret equals
$\tfrac12|R''(d^*)|(\hat d-d^*)^2+o((\hat d-d^*)^2)$.
\item $\hat d=\dlo$ if and only if
$(1-\dlo)\,\hat b\,r(\hat a+\hat b\,\dlo)\le1$, where $r=\phi/\Phi$. This
holds for sufficiently small positive $\hat b$, so the recommendation reaches
the floor before the fitted slope changes sign, and regret is then constant.
\end{enumerate}
\end{proposition}
\begin{proof}[Proof sketch]
(1) is the log-concavity argument above. (2) is a second-order expansion at an
interior maximum. (3) $\log\{(1-d)\Phi(\hat a+\hat b d)\}$ is concave, so the
maximizer is $\dlo$ exactly when its derivative there is non-positive.
\end{proof}

\begin{lemma}[Direction of the error]
\label{lem:direction}
Fix the value of the believed curve at $m$ and vary its slope $\hat b$. At an
interior optimum with $z=\hat a+\hat b\hat d$,
\[
\operatorname{sign}\frac{\partial\hat d}{\partial\hat b}
=\operatorname{sign}\bigl(1-\hat b(\hat d-m)(z+r(z))\bigr).
\]
\end{lemma}
\begin{proof}
Differentiate the first-order condition $(1-d)\hat b\,r(z)=1$ with respect to
$\hat b$ with $z=c+\hat b(d-m)$, using $r'(z)=-r(z)(z+r(z))$; the derivative
of the left side in $d$ is negative at a maximum.
\end{proof}

When the optimum is at or below the mean historical discount, a flatter curve
lowers the recommendation and a steeper one raises it. When the optimum is
well above the mean the directions reverse: a steeper curve makes the model
believe the deal is already won just past the pivot. So ``rescue confounding
causes under-discounting'' is not a general statement. Together with
Proposition~\ref{prop:regret} this gives the practical picture: modest biases
in the slope cost little, and the damage is concentrated where the
recommendation is pushed to a boundary.

\section{Correcting with Rep Leniency, and Why It Can Fail Silently}
\label{sec:instrument}

Suppose deals are handled by reps who differ in how much they discount for
reasons unrelated to any one deal. Extend \eqref{eq:first} and
\eqref{eq:outcome} to
\begin{equation}
d = f(x) + b_r + v, \qquad
Y=\ind\{h(x)+g(d)+\delta s_r+\beta q+\varepsilon>0\},
\label{eq:rep}
\end{equation}
where $b_r$ is the leniency and $s_r$ the selling skill of rep $r$, $x$ are
covariates, and $(v,q,\varepsilon)$ are independent of $(x,r)$ but may depend
on each other. If $\delta=0$ (exclusion) and reps are assigned independently
of $(v,q,\varepsilon)$, leniency is a valid instrument. We estimate it by the
leave-one-out mean, over a rep's other deals, of the discount residualized on
covariates \citep{dobbie2018effects}; the first-stage residual $\hat v$ enters
the win model as a control function \citep{blundell2004endogeneity}; and the
causal curve is the fitted probability averaged over the distribution of
$\hat v$.

\begin{proposition}[Leniency-aligned skill is undetectable]
\label{prop:undetectable}
In model \eqref{eq:rep}:
\begin{enumerate}
\item If $s_r=\kappa b_r$, the same data are generated, deal by deal, by
\[
Y=\ind\{\tilde h(x)+\tilde g(d)+\tilde u>0\},\quad
\tilde g(d)=g(d)+\delta\kappa d,\quad
\tilde h(x)=h(x)-\delta\kappa f(x),\quad
\tilde u=\beta q-\delta\kappa v+\varepsilon,
\]
in which there is no rep skill and $\tilde u$ is independent of $(x,r)$.
\item The two models have different causal responses to discount,
$g'(d)$ and $g'(d)+\delta\kappa$.
\item If $s_r=s(b_r)$ for any function $s$, then $Y$ is independent of $r$
given $(x,d,v)$.
\end{enumerate}
\end{proposition}
\begin{proof}
For (1), substitute $b_r=d-f(x)-v$ into $\delta\kappa b_r$ and collect terms.
(2) is immediate. For (3), $b_r=d-f(x)-v$ is a function of $(x,d,v)$, so
conditioning on $(x,d,v)$ fixes $s_r$, and the remaining randomness in $Y$ is
$(q,\varepsilon)$ given $v$, which does not depend on $r$.
\end{proof}

The second model in part 1 satisfies every assumption of the instrument, so
no test based on covariates, rep identity, discount and outcome can
distinguish it from the first, while a control function fitted to the data
recovers $\tilde g$. Part 3 says that tests for residual differences between
reps have no power against skill that is a function of leniency. Writing
$s_r=\kappa b_r+e_r$, only the component $e_r$ unrelated to leniency is
visible.

\begin{remark}[Sorting]
If managers route deals to reps on unrecorded quality, $\E[q\mid r]$ varies
with $b_r$ and enters the outcome equation exactly as $\delta s_r$ does. Our
experiments treat this case separately because it is the more plausible one in
sales organizations.
\end{remark}

\begin{remark}[A use for monotonicity]
When $\delta\kappa<0$ (lenient reps are weaker) and $g$ saturates, $\tilde g$
is eventually decreasing: the model without rep skill implies that large
discounts lose deals. An analyst who believes discounts never hurt can reject
it. So the monotonicity belief that is harmful as a constraint on the
uncorrected model is informative as a falsification check on the corrected
one. It gives no leverage when $\delta\kappa>0$.
\end{remark}

\paragraph{A sensitivity range.} Following \citet{conley2012plausibly}, we
refit the control function assuming leniency has a direct effect $\delta$ on
the index, for $\delta$ in a range $[-\Delta,\Delta]$, and recommend the
discount that minimizes the maximum regret across the fitted models. The
range is either fixed by the analyst or calibrated to the data by setting
$\Delta=\hat\tau$, the estimated standard deviation of residual rep effects:
the assumption is that hidden, leniency-aligned skill is no larger than the
skill differences that are visible.

\section{Benchmark}
\label{sec:benchmark}

The benchmark generates books of closed deals with a known causal response.
Each deal has a segment, region, industry, list price, competitor indicator
and quarter-end indicator, drawn as in the synthetic generator of a
commercial deal-analytics product; a rep; and a latent quality. The discount
is the sum of a segment base level, competitor and quarter-end effects, the
rep's bias, and a residual with standard deviation 3.5 points whose
correlation with latent quality is the dial $\rho$. The win probability is
logistic in a covariate index, the latent quality, and a saturating effect of
discount, $1.3\tanh(d/0.05)$, which is essentially exhausted by a 10\%
discount. The spread of discounts is the same at every $\rho$, so only the
reason for discounting changes.

Further dials are the weight of latent quality in the outcome, the number of
training deals, the number of reps, a rep skill effect and its correlation
with leniency, and the rule assigning reps to deals: at random, by segment
(territories), or by a noisy reading of latent quality (sorting).

Because the response is known, the causal curve for each deal is computed by
integrating over latent quality, and each recommendation is scored by regret:
the revenue lost against the best discount for that deal, as a percentage of
the revenue at that best discount, weighted by list price. We report regret
for an unrestricted sweep over discounts from 0 to 40\% and, where stated, for
a sweep clipped to the 5th to 95th percentile of the discounts observed in the
deal's segment, a guardrail used in practice. All results are means over 30
independently generated books with standard errors, except the capacity sweep
(10 books). Recommendations are evaluated on fresh deals.

\todo{State licence and repository. Add the real-covariate check on public
auto-lending data if the dataset can be obtained.}

\section{Experiments}
\label{sec:experiments}

\subsection{The monotone constraint}

We train gradient-boosted trees \todo{cite LightGBM} with the configuration
of a production win model (50 trees, 7 leaves, at least 150 deals per leaf),
with and without a monotone constraint on discount. Table~\ref{tab:constraint}
reports regret on books of 20{,}000 deals.

\begin{table}[t]
\centering\small
\begin{tabular}{rrrrrr}
\toprule
$\rho$ & Constrained & Unconstrained & Difference & Constrained (clip) & Difference (clip) \\
\midrule
$-0.9$ & 30.5 (1.0) & 19.5 (1.2) & $+11.0$ (1.0) & 5.2 (0.3) & $+0.9$ (0.1) \\
$-0.6$ & 22.1 (0.8) & 13.9 (0.8) & $+8.2$ (0.6) & 3.9 (0.2) & $+0.7$ (0.1) \\
$-0.3$ & 15.9 (0.7) & 11.7 (0.6) & $+4.3$ (0.5) & 3.0 (0.2) & $+0.3$ (0.1) \\
0 & 9.3 (0.6) & 8.7 (0.5) & $+0.6$ (0.3) & 1.8 (0.1) & $0.0$ (0.0) \\
$+0.3$ & 4.2 (0.4) & 4.3 (0.3) & $-0.1$ (0.3) & 1.0 (0.1) & $-0.1$ (0.0) \\
$+0.6$ & 1.8 (0.2) & 2.2 (0.2) & $-0.4$ (0.2) & 0.6 (0.1) & $-0.1$ (0.1) \\
$+0.9$ & 1.4 (0.2) & 2.1 (0.3) & $-0.7$ (0.2) & 0.7 (0.1) & $-0.4$ (0.1) \\
\bottomrule
\end{tabular}

\caption{Regret (\% of optimal revenue; mean and standard error over 30
books) of a tree win model with and without a monotone constraint on
discount, 20{,}000 training deals. ``Difference'' is constrained minus
unconstrained, paired by book. ``Clip'' restricts recommendations to the
segment's observed discount range.}
\label{tab:constraint}
\end{table}

Under rescue confounding the constraint raises regret by 4.3 to 11.0 points on
the unrestricted sweep, as Proposition~\ref{prop:saturating} predicts. Under
whale confounding it makes little difference. The constraint does what it was
introduced for: the unconstrained model produces a visibly decreasing curve
for most deals under rescue confounding and the constrained model for none.
Clipping recommendations to the observed range reduces the worst-case regret
of the constrained model from 30.5\% to 5.2\% and the cost of the constraint
to under a point. In this configuration it is the clip, not the constraint,
that protects the decision. Results at 2{,}000 deals are similar
(Appendix~\ref{app:tables}).

\subsection{Model capacity}

Table~\ref{tab:capacity} varies the capacity of the tree model.

\begin{table}[t]
\centering\small
\begin{tabular}{lrrrrrr}
\toprule
Capacity & $\rho=-0.6$ & $\rho=0$ & $\rho=+0.6$ & Diff., $-0.6$ & Diff., $+0.6$ & Mean rec. \\
\midrule
50 trees, 7 leaves, 150 per leaf (shipped) & 22.6 (1.2) & 9.0 (0.7) & 1.4 (0.3) & $+7.9$ (0.8) & $-0.5$ (0.2) & 5.0\% \\
200 trees, 7 leaves, 150 per leaf & 11.9 (0.8) & 1.8 (0.4) & 6.7 (0.5) & $+4.2$ (0.5) & $+0.3$ (0.2) & 7.3\% \\
500 trees, 7 leaves, 50 per leaf & 11.0 (0.9) & 1.8 (0.2) & 11.0 (0.7) & $+3.0$ (1.3) & $+1.0$ (0.4) & 8.2\% \\
500 trees, 15 leaves, 50 per leaf & 10.3 (1.0) & 2.7 (0.3) & 12.2 (0.6) & $+0.1$ (2.1) & $+2.4$ (0.3) & 8.9\% \\
500 trees, 15 leaves, 20 per leaf, rate 0.1 & 11.9 (0.9) & 4.9 (0.7) & 12.8 (0.6) & $-1.5$ (1.9) & $+2.3$ (0.3) & 10.3\% \\
\bottomrule
\end{tabular}

\caption{Regret of the constrained tree model at five capacities, 20{,}000
training deals (10 books). ``Diff.'' is constrained minus unconstrained.
``Mean rec.'' is the mean recommended discount at $\rho=0$; the mean optimal
discount is 8.6\%.}
\label{tab:capacity}
\end{table}

Three things stand out. First, with no confounding at all the production
configuration loses 9.0\% of revenue, because fifty boosting rounds at a small
learning rate under-read the steep initial part of the response; two hundred
rounds lose 1.8\%. Second, added capacity moves the exposure from one
direction of confounding to the other: the small model is nearly unaffected by
whale confounding and badly affected by rescue confounding, and at 500 trees
the reverse holds. The small model's under-reading happens to offset the
upward bias from whale confounding. Third, the cost of the constraint under
rescue confounding shrinks as capacity grows and is within noise at the two
largest settings, while a cost under whale confounding appears, which
Proposition~\ref{prop:saturating} does not predict and which we attribute to
finite samples: the constraint removes the noisy dips that otherwise stop the
model from recommending still more discount. \todo{Rerun with 30 books and
with the clipped sweep.}

\subsection{The leniency control function}

To separate the correction from the capacity of the tree model, we use a
logistic win model with a cubic spline in discount, fitted without correction
(``naive''), with the leniency control function (``corrected''), and with the
latent quality observed (``oracle'').

\begin{table}[t]
\centering\small
\begin{tabular}{rrrrrrr}
\toprule
$\rho$ & Naive & Corrected & Oracle & Naive & Corrected & Oracle \\
\midrule
$-0.9$ & 6.0 (0.5) & 1.2 (0.4) & 0.3 (0.1) & 8.6 (1.5) & 7.4 (1.5) & 7.0 (1.5) \\
$-0.6$ & 2.3 (0.2) & 0.5 (0.1) & 0.2 (0.0) & 4.2 (0.8) & 8.1 (1.4) & 6.9 (1.3) \\
$-0.3$ & 0.6 (0.1) & 0.3 (0.1) & 0.2 (0.0) & 3.6 (0.9) & 6.5 (1.2) & 5.3 (1.2) \\
0 & 0.2 (0.0) & 0.3 (0.0) & 0.2 (0.0) & 4.2 (1.1) & 4.3 (1.0) & 4.8 (1.1) \\
$+0.3$ & 4.7 (0.6) & 0.3 (0.0) & 0.2 (0.1) & 11.4 (1.1) & 4.1 (1.0) & 3.8 (1.0) \\
$+0.6$ & 13.7 (0.7) & 0.4 (0.1) & 0.3 (0.1) & 17.2 (0.8) & 3.3 (0.8) & 3.1 (0.8) \\
$+0.9$ & 18.4 (0.5) & 0.5 (0.1) & 0.3 (0.0) & 18.8 (0.5) & 3.3 (0.8) & 3.8 (1.0) \\
\bottomrule
\end{tabular}

\caption{Regret of a spline win model, unrestricted sweep. Columns 2 to 4:
20{,}000 training deals. Columns 5 to 7: 2{,}000.}
\label{tab:spline}
\end{table}

\begin{figure}[t]
\centering
\begin{tikzpicture}
\begin{axis}[width=0.78\textwidth, height=0.46\textwidth,
  xlabel={Confounding $\rho$ (rescue $<0<$ whale)},
  ylabel={Regret (\% of optimal revenue)},
  xmin=-1, xmax=1, ymin=0, ymax=34,
  legend pos=north east, legend cell align=left]
\addplot[black, thick, mark=square*, error bars/.cd, y dir=both, y explicit]
  table[x=rho, y=treecon, y error=treeconse] {data/regret.dat};
\addlegendentry{Tree, constrained}
\addplot[black, thick, dashed, mark=square, error bars/.cd, y dir=both, y explicit]
  table[x=rho, y=treeunc, y error=treeuncse] {data/regret.dat};
\addlegendentry{Tree, unconstrained}
\addplot[black, thick, dotted, mark=*, error bars/.cd, y dir=both, y explicit]
  table[x=rho, y=splinenaive, y error=splinenaivese] {data/regret.dat};
\addlegendentry{Spline, naive}
\addplot[black, thick, dashdotted, mark=o, error bars/.cd, y dir=both, y explicit]
  table[x=rho, y=splinecf, y error=splinecfse] {data/regret.dat};
\addlegendentry{Spline, corrected}
\end{axis}
\end{tikzpicture}
\caption{Regret against the strength and direction of confounding, 20{,}000
training deals, unrestricted sweep. Bars are one standard error over 30
books. The tree model at production capacity is exposed to rescue
confounding and the spline model to whale confounding.}
\label{fig:regret}
\end{figure}

At 20{,}000 deals the corrected model is within a point of the oracle at
every setting and reduces the worst regret of the naive model from 18.4\% to
0.5\% (Table~\ref{tab:spline}, Figure~\ref{fig:regret}). The naive spline
model is harmed mainly by whale confounding, the opposite of the production
tree model.

At 2{,}000 deals the corrected model again tracks the oracle, and both lose 3
to 8\%. Under rescue confounding the naive model does better than either,
because its bias offsets a small-sample tendency of the spline model to
over-discount. The correction is doing its job; at this sample size the job is
not worth doing in that direction.

Table~\ref{tab:reps} varies the number of reps. Six reps do about as well as
fifty: with random assignment the number of deals, not the number of reps,
limits the correction.

\begin{table}[t]
\centering\small
\begin{tabular}{rrrrrrr}
\toprule
Reps & $-0.6$ & $0$ & $+0.6$ & $-0.6$ & $0$ & $+0.6$ \\
\midrule
6 & $-0.8$ (0.7) & $-0.0$ (0.0) & $-14.1$ (0.7) & $+4.8$ (1.5) & $+0.5$ (0.8) & $-9.8$ (0.9) \\
12 & $+0.2$ (0.9) & $+0.3$ (0.2) & $-13.2$ (0.9) & $+0.7$ (1.2) & $+0.3$ (0.7) & $-10.3$ (0.9) \\
25 & $-0.1$ (0.7) & $+0.1$ (0.1) & $-12.6$ (0.6) & $+4.0$ (1.0) & $+0.7$ (0.6) & $-10.8$ (0.9) \\
50 & $-1.0$ (0.4) & $+0.1$ (0.0) & $-11.0$ (0.5) & $+2.2$ (1.1) & $+0.0$ (0.5) & $-11.5$ (0.8) \\
\bottomrule
\end{tabular}

\caption{Corrected minus naive regret by number of reps. Columns 2 to 4:
20{,}000 training deals at $\rho=-0.6,0,+0.6$. Columns 5 to 7: 2{,}000.}
\label{tab:reps}
\end{table}

\subsection{Rep skill}

Table~\ref{tab:exclusion} adds rep skill worth 0.3 on the logit scale per
standard deviation, either exactly aligned with leniency (lenient reps are
weaker), partly aligned, or unrelated to it.

\begin{table}[t]
\centering\small
\begin{tabular}{lrrrrrr}
\toprule
Scenario & $\rho$ & Naive & Corrected & Robust (calibrated) & Robust (fixed) & Segment median \\
\midrule
valid & $-0.6$ & 2.2 (0.2) & 2.9 (0.9) & 2.9 (0.9) & 4.0 (0.1) & 6.9 (0.2) \\
 & 0 & 0.2 (0.1) & 1.3 (0.5) & 1.3 (0.5) & 3.7 (0.1) & 6.8 (0.2) \\
 & $+0.6$ & 13.8 (0.9) & 1.4 (0.7) & 1.4 (0.7) & 3.4 (0.2) & 6.8 (0.1) \\
\midrule
aligned & $-0.6$ & 9.5 (0.5) & 9.5 (0.9) & 9.4 (0.9) & 3.2 (0.2) & 6.9 (0.2) \\
 & 0 & 2.1 (0.2) & 8.0 (0.6) & 7.9 (0.6) & 3.1 (0.2) & 6.9 (0.2) \\
 & $+0.6$ & 1.4 (0.5) & 7.6 (0.5) & 7.6 (0.5) & 3.6 (0.2) & 6.9 (0.1) \\
\midrule
partial & $-0.6$ & 8.6 (0.5) & 8.2 (0.9) & 2.8 (0.2) & 2.3 (0.2) & 7.0 (0.2) \\
 & 0 & 1.8 (0.2) & 7.0 (0.6) & 2.9 (0.2) & 2.3 (0.1) & 7.0 (0.2) \\
 & $+0.6$ & 1.9 (0.7) & 6.6 (0.5) & 3.2 (0.2) & 2.7 (0.2) & 7.0 (0.1) \\
\midrule
unaligned & $-0.6$ & 2.3 (0.2) & 2.8 (0.7) & 4.0 (0.1) & 4.0 (0.1) & 6.9 (0.2) \\
 & 0 & 0.2 (0.0) & 0.8 (0.3) & 3.7 (0.1) & 3.8 (0.1) & 6.9 (0.2) \\
 & $+0.6$ & 14.3 (0.9) & 0.6 (0.2) & 3.4 (0.2) & 3.5 (0.2) & 6.9 (0.1) \\
\midrule
worst case &  & 14.3 & 9.5 & 9.4 & 4.0 & 7.0 \\
\bottomrule
\end{tabular}

\caption{Regret when reps differ in skill, 20{,}000 training deals. ``Robust''
is the minimax-regret recommendation over a sensitivity range that is
calibrated to residual rep differences or fixed at $\pm0.3$. ``Segment
median'' recommends the segment's median historical discount.}
\label{tab:exclusion}
\end{table}

With aligned skill the corrected model loses 7.6 to 9.5\%, and with no
confounding or whale confounding it is far worse than the naive model. The
test for residual rep differences flags these books no more often than books
with a valid instrument (Table~\ref{tab:reptest}), as
Proposition~\ref{prop:undetectable} requires. It flags skill that is not
aligned with leniency in essentially every book and recovers its magnitude.

\begin{table}[t]
\centering\small
\begin{tabular}{lrrrr}
\toprule
Scenario & $\hat\tau$ & Flagged & $\hat\tau$ & Flagged \\
\midrule
valid & 0.01 & 6\% & 0.03 & 7\% \\
aligned & 0.01 & 9\% & 0.03 & 4\% \\
partial & 0.28 & 100\% & 0.27 & 98\% \\
unaligned & 0.29 & 100\% & 0.29 & 100\% \\
\bottomrule
\end{tabular}

\caption{The rep test: mean estimated spread of residual rep effects
($\hat\tau$, logit scale) and share of books flagged at the 5\% level.
Columns 2 and 3: 20{,}000 training deals. Columns 4 and 5: 2{,}000. The true
visible spread is 0, 0, 0.29 and 0.30.}
\label{tab:reptest}
\end{table}

The calibrated range costs nothing when the instrument is valid and protects
when some skill is visible. It cannot protect in the fully aligned case. The
fixed range has the best worst case, 4.0\%, at a price of about three points
when nothing is wrong; it contains the true direct effect by construction in
these scenarios. An alternative design, abstaining on deals whose best
discount moves by more than five points across the range and falling back to
the segment median, was worse than the robust recommendation it replaced by
3.2 to 4.4 points, because the fallback is itself a poor recommendation.

\subsection{Non-random assignment}

Finally we let assignment depend on the deal (Table~\ref{tab:assignment}).

\begin{table}[t]
\centering\small
\begin{tabular}{lrrrrrr}
\toprule
Assignment & $\rho$ & Naive & Corrected & Robust (fixed) & Balance & Rep test \\
\midrule
Random & $-0.6$ & 2.2 (0.2) & 2.9 (0.9) & 4.0 (0.1) & 3\% & 6\% \\
 & 0 & 0.2 (0.1) & 1.3 (0.5) & 3.7 (0.1) &  &  \\
 & $+0.6$ & 13.8 (0.9) & 1.4 (0.7) & 3.4 (0.2) &  &  \\
\midrule
Territory & $-0.6$ & 2.3 (0.4) & 4.1 (1.3) & 4.6 (0.3) & 100\% & 1\% \\
 & 0 & 0.8 (0.4) & 3.4 (1.2) & 4.2 (0.3) &  &  \\
 & $+0.6$ & 16.5 (0.7) & 4.0 (1.3) & 5.0 (0.8) &  &  \\
\midrule
Sorting $+0.3$ & $-0.6$ & 1.2 (0.2) & 18.6 (0.8) & 7.4 (0.3) & 5\% & 29\% \\
 & 0 & 5.3 (0.8) & 15.7 (0.9) & 7.1 (0.2) &  &  \\
 & $+0.6$ & 17.1 (0.6) & 12.4 (1.0) & 6.7 (0.3) &  &  \\
\midrule
Sorting $+0.6$ & $-0.6$ & 1.4 (0.3) & 21.2 (0.3) & 13.7 (0.6) & 5\% & 79\% \\
 & 0 & 14.2 (0.6) & 21.1 (0.3) & 13.3 (0.6) &  &  \\
 & $+0.6$ & 18.8 (0.4) & 19.5 (0.5) & 11.7 (0.5) &  &  \\
\midrule
Sorting $-0.3$ & $-0.6$ & 3.8 (0.3) & 1.6 (0.2) & 0.8 (0.1) & 7\% & 44\% \\
 & 0 & 0.9 (0.1) & 2.8 (0.3) & 1.0 (0.1) &  &  \\
 & $+0.6$ & 7.3 (1.3) & 3.8 (0.3) & 1.2 (0.2) &  &  \\
\midrule
Sorting $-0.6$ & $-0.6$ & 5.2 (0.3) & 5.8 (0.5) & 3.0 (0.2) & 7\% & 86\% \\
 & 0 & 2.5 (0.2) & 9.4 (0.7) & 4.2 (0.2) &  &  \\
 & $+0.6$ & 2.2 (0.7) & 15.2 (1.4) & 4.5 (0.3) &  &  \\
\bottomrule
\end{tabular}

\caption{Regret under non-random assignment of reps, 20{,}000 training deals.
``Territory'': each deal goes to a rep whose home segment matches.
``Sorting $\pm c$'': the correlation between a deal's latent quality and the
leniency rank of its rep is $c$ in magnitude; positive sends strong deals to
lenient reps. ``Balance'' and ``Rep test'' are the share of books flagged at
the 5\% level.}
\label{tab:assignment}
\end{table}

Assignment by territory is flagged by the balance test in every book and is
nearly harmless once segment is controlled for. Sorting on latent quality is
not flagged by the balance test, since the book does not record what the
sorting is based on, and it makes the corrected model lose 12 to 21\% when
strong deals go to lenient reps. The first-stage $F$ statistic is, if
anything, larger under sorting. The rep test has some power on large books,
which comes from the mismatch between rank-based sorting and the additive
form of the win model and should not be relied on. In nine of the twelve
sorting cells the naive model has lower regret than the corrected one.

\section{Guidance}
\label{sec:guidance}

\todo{Turn into a table indexed by book size, model capacity and what is
known about assignment.} The results suggest the following.

\begin{itemize}
\item Check capacity before confounding. A win model that under-fits the
response is biased toward too little discount whatever the data.
\item Do not expect a monotone constraint to fix the recommendation. Restrict
recommendations to the range of discounts the cohort has actually received.
\item Use a leniency instrument only where assignment follows a known rule,
such as round-robin routing or territories recorded in the data. A balance
test is necessary and not sufficient.
\item If the corrected curve says large discounts lose deals, treat that as
evidence against the instrument.
\item When assignment is doubtful, prefer a recommendation that is robust over
a stated range of direct effects to either trusting the correction or falling
back to historical practice.
\item On small books, do not correct.
\end{itemize}

\section{Limitations}
\label{sec:limits}

All outcomes are simulated. The benchmark is one family of generators with a
saturating response, normally distributed latent quality, and an additive
discount equation; the propositions are population statements for a single
cohort. The spline win model is close to the generator's functional form. We
tested one mechanism each for skill and for sorting. The paper says nothing
about which strength or direction of confounding real deal books exhibit;
\citet{phillips2015field} is the one estimate we know of, in auto lending.

\todo{Conflict of interest statement. Statement on the use of AI tools.}

\appendix
\section{Additional Tables}
\label{app:tables}

\begin{table}[h]
\centering\small
\begin{tabular}{rrrrrr}
\toprule
$\rho$ & Constrained & Unconstrained & Difference & Constrained (clip) & Difference (clip) \\
\midrule
$-0.9$ & 34.8 (1.8) & 24.7 (2.5) & $+10.1$ (1.8) & 5.3 (0.3) & $+1.3$ (0.2) \\
$-0.6$ & 27.7 (1.8) & 19.9 (2.1) & $+7.8$ (1.3) & 4.2 (0.3) & $+0.7$ (0.2) \\
$-0.3$ & 21.5 (1.8) & 16.4 (2.0) & $+5.2$ (0.9) & 3.1 (0.3) & $+0.4$ (0.1) \\
0 & 14.0 (1.4) & 11.9 (1.5) & $+2.2$ (0.7) & 2.5 (0.3) & $+0.1$ (0.1) \\
$+0.3$ & 9.8 (1.1) & 8.8 (1.1) & $+0.9$ (0.4) & 2.1 (0.2) & $+0.1$ (0.1) \\
$+0.6$ & 7.7 (0.9) & 7.5 (0.9) & $+0.2$ (0.2) & 2.1 (0.2) & $-0.0$ (0.1) \\
$+0.9$ & 5.4 (0.7) & 5.8 (0.7) & $-0.4$ (0.2) & 2.7 (0.3) & $+0.0$ (0.1) \\
\bottomrule
\end{tabular}

\caption{As Table~\ref{tab:constraint}, with 2{,}000 training deals.}
\end{table}

\section{Reproducibility Checklist}
\todo{Insert the completed JAIR reproducibility checklist.}

\bibliographystyle{plainnat}
\bibliography{refs}
\end{document}
